% !TEX spellcheck = en_US



% ~~~~~~~~~~~~~~~~~~~~~~~~~~~~~~~~~~~~~~~ V
% (NC) 2026 Haluk Bingol
% github.com/halukbingol/LaTeX-Templates
%
% Licensees may copy, distribute, display, and perform the work and make derivative works 
% and remixes based on it only for non-commercial purposes. 
% ~~~~~~~~~~~~~~~~~~~~~~~~~~~~~~~~~~~~~~~ A




% v20261030T1200

\documentclass[aspectratio=169,10pt,compress]{beamer}




% ~~~~~~~~~~~~~~~~~~~~~~~~~~~~~~~~~~~~~~~~~~~~~~~~~~~~~~~~~~~~~~~~~ V hblib
	\usepackage{../../hblib/hblibPresentation} % lib presentation: theme, i/N, \hSection, algorithm2e, ...
	\usepackage{../../hblib/hblibCodeoutput} % lib code + output: \lstcode, \lstterminal, ...
	\usepackage{../../hblib/hblibDatastructures} % lib data structures: \hString, \hStack, \hQueue, \hArray, \hGrid
% ~~~~~~~~~~~~~~~~~~~~~~~~~~~~~~~~~~~~~~~~~~~~~~~~~~~~~~~~~~~~~~~~~ A hblib




% +++++++++++++++++++++++++++++++++++++++ V specific
\DontPrintSemicolon % in algorithms
% +++++++++++++++++++++++++++++++++++++++ A specific



\title[Arrays ND: Freshman]{
	Multi-Dimensional Arrays\\
	Freshman
}
\subtitle{
	\cisOOO\ -- Level 1: first steps and common errors
}
\author[Erkovan]{Sait Erkovan}

\institute[FCIS]{
	Faculty of Computer and Information Sciences\\
	Yeditepe University
}
\date{{\tiny \hbTimeStamp}}

\begin{document}




% ~~~~~~~~~~~~~~~~~~~~~~~~~~~~~~~~~~~~~~ V frm
\begin{frame}[t,fragile]
	\titlepage
\end{frame}
% ~~~~~~~~~~~~~~~~~~~~~~~~~~~~~~~~~~~~~~ A frm




% ~~~~~~~~~~~~~~~~~~~~~~~~~~~~~~~~~~~~~~ V frm
\begin{frame}[t,fragile] \frametitle{Outline}
	\begin{columns}[T]
		\begin{column}{0.48\textwidth}
			\tableofcontents[sections={1-3}, hideallsubsections]
		\end{column}
		\begin{column}{0.48\textwidth}
			\tableofcontents[sections={4-7}, hideallsubsections]
		\end{column}
	\end{columns}

	\vspace{1em}
	\small It builds on \hCode{array1D-1-freshman}: the 1D syntax and its errors are not repeated here.
\end{frame}
% ~~~~~~~~~~~~~~~~~~~~~~~~~~~~~~~~~~~~~~ A frm




% ~~~~~~~~~~~~~~~~~~~~~~~~~~~~~~~~~~~~~~ V
\hSection[Create]{Creating a table}
% ~~~~~~~~~~~~~~~~~~~~~~~~~~~~~~~~~~~~~~ A




% ~~~~~~~~~~~~~~~~~~~~~~~~~~~~~~~~~~~~~~ V frm
\begin{frame}[t,fragile] \frametitle{Three ways to create a table}


% +++++++++++++++++++++++++++++++++++++++ V lst
%: -Lst. three ways
\begin{lstlisting}[style=codesnippet, language=Java]
int[][] a = new int[3][4];               // 1. sizes: 3 rows, 4 columns, all 0

int[][] b = {{72, 85, 90},               // 2. values: one pair of braces per row
             {55, 61, 70}};

int[][] c = new int[3][];                // 3. rows later: 3 rows, none created yet
c[0] = new int[2];
c[1] = new int[5];                       //    rows may differ in length
c[2] = new int[1];
\end{lstlisting}
% +++++++++++++++++++++++++++++++++++++++ A lst


	\begin{itemize}\small
		\item \hCode{[rows][columns]}: the first number is always the rows
		\item Way 2 is only for a declaration: \hCode{b = \{\{1\}\};} alone does not compile, as in 1D
		\item Way 3 makes a \hDefined{ragged} table; until a row is created, it does not exist
	\end{itemize}
\end{frame}
% ~~~~~~~~~~~~~~~~~~~~~~~~~~~~~~~~~~~~~~ A frm




% ~~~~~~~~~~~~~~~~~~~~~~~~~~~~~~~~~~~~~~ V frm
\begin{frame}[t,fragile] \frametitle{A first program}
	\lstcoderightstep[0.66][basicstyle=\ttfamily\scriptsize]{codeoutput/ScoresTable.java}{Java}{codeoutput/ScoresTable.txt}

	\footnotesize Named sizes (\hCode{STUDENTS}, \hCode{QUIZZES}) say which number is which, and change in one place.
	Unwritten cells keep the start value 0.
\end{frame}
% ~~~~~~~~~~~~~~~~~~~~~~~~~~~~~~~~~~~~~~ A frm




% ~~~~~~~~~~~~~~~~~~~~~~~~~~~~~~~~~~~~~~ V frm
\begin{frame}[t,fragile] \frametitle{Initializer lists: one pair of braces per row}
	\lstcoderightstep[0.66][basicstyle=\ttfamily\scriptsize]{codeoutput/InitRows.java}{Java}{codeoutput/InitRows.txt}

	\footnotesize The outer braces hold the rows; each inner pair is one row. \hCode{\{\}} is an empty row of length 0.
\end{frame}
% ~~~~~~~~~~~~~~~~~~~~~~~~~~~~~~~~~~~~~~ A frm




% ~~~~~~~~~~~~~~~~~~~~~~~~~~~~~~~~~~~~~~ V frm
\begin{frame}[t,fragile] \frametitle{Reading and writing cells}


% +++++++++++++++++++++++++++++++++++++++ V tbl
%: -Tbl. cell syntax
\begin{table}
		\small
		\begin{tabular}{@{}L{0.36\textwidth}L{0.56\textwidth}@{}}
			\toprule
			\textbf{Write} & \textbf{Meaning} \\
			\midrule
			\hCode{t[r][c]}                        & the cell in row \hCode{r}, column \hCode{c}: read it \\
			\hCode{t[r][c] = 7;} \ \hCode{t[r][c]++;} & write it, change it, like any variable \\
			\hCode{t[r]}                           & row \hCode{r}: an \hCode{int[]} \\
			\hCode{t.length}                       & the number of rows \\
			\hCode{t[r].length}                    & the number of cells in row \hCode{r} \\
			\hCode{t[t.length - 1]}                & the last row \\
			\hCode{t[r][t[r].length - 1]}          & the last cell of row \hCode{r} \\
			\bottomrule
		\end{tabular}
\end{table}%
% +++++++++++++++++++++++++++++++++++++++ A tbl


	\vspace{0.3em}
	\small There is no ``number of columns'' in Java: only the length of each row.
\end{frame}
% ~~~~~~~~~~~~~~~~~~~~~~~~~~~~~~~~~~~~~~ A frm




% ~~~~~~~~~~~~~~~~~~~~~~~~~~~~~~~~~~~~~~ V
\hSection[Loops]{Loops over a table}
% ~~~~~~~~~~~~~~~~~~~~~~~~~~~~~~~~~~~~~~ A




% ~~~~~~~~~~~~~~~~~~~~~~~~~~~~~~~~~~~~~~ V frm
\begin{frame}[t,fragile] \frametitle{Index loops and the enhanced for}
	\lstcoderightstep[0.66][basicstyle=\ttfamily\scriptsize]{codeoutput/LoopKinds.java}{Java}{codeoutput/LoopKinds.txt}

	\footnotesize Lines 5--10: index loops, with \hCode{w} and \hCode{d}. Lines 12--16: the enhanced for, row by row, cell by cell.
\end{frame}
% ~~~~~~~~~~~~~~~~~~~~~~~~~~~~~~~~~~~~~~ A frm




% ~~~~~~~~~~~~~~~~~~~~~~~~~~~~~~~~~~~~~~ V frm
\begin{frame}[t,fragile] \frametitle{Which loop?}


% +++++++++++++++++++++++++++++++++++++++ V tbl
%: -Tbl. which loop
\begin{table}
		\small
		\begin{tabular}{@{}L{0.48\textwidth}cc@{}}
			\toprule
			\textbf{You need to \dots} & \textbf{index loops} & \textbf{enhanced for} \\
			\midrule
			read every cell, row by row                 & yes & \hEmph{yes, simpler} \\
			change cells                                & yes & no \\
			know the row and the column (print them)    & yes & no \\
			go column by column                         & yes & no \\
			visit a part: one row, one column, a diagonal & yes & only one row \\
			\bottomrule
		\end{tabular}
\end{table}%
% +++++++++++++++++++++++++++++++++++++++ A tbl




% +++++++++++++++++++++++++++++++++++++++ V lst
%: -Lst. both loops
\begin{lstlisting}[style=codesnippet, language=Java]
for (int r = 0; r < t.length; r++) {
    for (int c = 0; c < t[r].length; c++) { ... t[r][c] ... }
}
for (int[] row : t) { for (int x : row) { ... x ... } }
\end{lstlisting}
% +++++++++++++++++++++++++++++++++++++++ A lst


\end{frame}
% ~~~~~~~~~~~~~~~~~~~~~~~~~~~~~~~~~~~~~~ A frm




% ~~~~~~~~~~~~~~~~~~~~~~~~~~~~~~~~~~~~~~ V
\hSection[Compile]{Errors the compiler finds}
% ~~~~~~~~~~~~~~~~~~~~~~~~~~~~~~~~~~~~~~ A




% ~~~~~~~~~~~~~~~~~~~~~~~~~~~~~~~~~~~~~~ V frm
\begin{frame}[t,fragile] \frametitle{Brackets: one pair per index}
	\lstoutput[basicstyle=\ttfamily\scriptsize]{codeoutput/CommaIndex.txt}

	\footnotesize \hCode{t[1, 2]} is not Java. Write \hCode{t[1][2]}.

	\vspace{0.4em}
	\lstoutput[basicstyle=\ttfamily\scriptsize]{codeoutput/NoRowCount.txt}

	\footnotesize The number of rows is required; the columns may come later: \hCode{new int[3][]}.
\end{frame}
% ~~~~~~~~~~~~~~~~~~~~~~~~~~~~~~~~~~~~~~ A frm




% ~~~~~~~~~~~~~~~~~~~~~~~~~~~~~~~~~~~~~~ V frm
\begin{frame}[t,fragile] \frametitle{One index too many, or too few}
	\lstoutput[basicstyle=\ttfamily\scriptsize]{codeoutput/TooManyIndices.txt}

	\begin{itemize}\footnotesize
		\item \hCode{t[0][0]} is already an \hCode{int}; a third \hCode{[0]} needs an array: ``array required''
		\item \hCode{t[0]} is a whole row, an \hCode{int[]}: it does not fit in an \hCode{int}
		\item Count the brackets: an \hCode{int[][]} needs two indices to give an \hCode{int}
	\end{itemize}
\end{frame}
% ~~~~~~~~~~~~~~~~~~~~~~~~~~~~~~~~~~~~~~ A frm




% ~~~~~~~~~~~~~~~~~~~~~~~~~~~~~~~~~~~~~~ V frm
\begin{frame}[t,fragile] \frametitle{The types must match}
	\lstoutput[basicstyle=\ttfamily\scriptsize]{codeoutput/WrongTypes.txt}

	\begin{itemize}\footnotesize
		\item Every cell has the element type: \hCode{3.5} cannot go into an \hCode{int} table
		\item \hCode{int[][]} and \hCode{int[]} are different types: a table is not a row
	\end{itemize}
\end{frame}
% ~~~~~~~~~~~~~~~~~~~~~~~~~~~~~~~~~~~~~~ A frm




% ~~~~~~~~~~~~~~~~~~~~~~~~~~~~~~~~~~~~~~ V frm
\begin{frame}[t,fragile] \frametitle{Reading these messages}


% +++++++++++++++++++++++++++++++++++++++ V tbl
%: -Tbl. messages
\begin{table}
		\footnotesize
		\begin{tabular}{@{}L{0.36\textwidth}L{0.30\textwidth}L{0.26\textwidth}@{}}
			\toprule
			\textbf{Message} & \textbf{Cause} & \textbf{Fix} \\
			\midrule
			\hCode{']' expected}                                & \hCode{t[1, 2]}; or \hCode{new int[][3]} & \hCode{t[1][2]}; \hCode{new int[3][]} \\
			\hCode{array required, but int found}               & too many indices                     & one index per \hCode{[]} \\
			\hCode{int[] cannot be converted to int}            & too few indices                      & add the column index \\
			\hCode{int[][] cannot be converted to int[]}        & a table where a row is expected      & pick a row: \hCode{t[r]} \\
			\hCode{possible lossy conversion from double to int} & a value of the wrong type           & the right element type \\
			\bottomrule
		\end{tabular}
\end{table}%
% +++++++++++++++++++++++++++++++++++++++ A tbl


	\vspace{0.3em}
	\small Read the line number and the caret \hCode{\^{}} first; the message says what was expected there.
\end{frame}
% ~~~~~~~~~~~~~~~~~~~~~~~~~~~~~~~~~~~~~~ A frm




% ~~~~~~~~~~~~~~~~~~~~~~~~~~~~~~~~~~~~~~ V
\hSection[Runtime]{Errors at run time}
% ~~~~~~~~~~~~~~~~~~~~~~~~~~~~~~~~~~~~~~ A




% ~~~~~~~~~~~~~~~~~~~~~~~~~~~~~~~~~~~~~~ V frm
\begin{frame}[t,fragile] \frametitle{Rows and columns swapped}
	\footnotesize \hQuestion{2 rows, 3 columns. Which index goes out of bounds?}

	\lstcodedownstep[basicstyle=\ttfamily\scriptsize, firstline=3, lastline=9, firstnumber=3][basicstyle=\ttfamily\fontsize{6.4}{7.6}\selectfont]{codeoutput/SwappedIndex.java}{Java}{codeoutput/SwappedIndex.txt}

	\footnotesize \uncover<2->{\hCode{r} counts to 3 but there are only 2 rows: \hCode{t[2]} fails at line 6.
	On a square table this bug stays hidden. Rule: \hCode{r} with \hCode{t.length}, \hCode{c} with \hCode{t[r].length}.}
\end{frame}
% ~~~~~~~~~~~~~~~~~~~~~~~~~~~~~~~~~~~~~~ A frm




% ~~~~~~~~~~~~~~~~~~~~~~~~~~~~~~~~~~~~~~ V frm
\begin{frame}[t,fragile] \frametitle{A row that was never created}
	\lstcodedownstep[basicstyle=\ttfamily\scriptsize][basicstyle=\ttfamily\fontsize{6.4}{7.6}\selectfont]{codeoutput/MissingRow.java}{Java}{codeoutput/MissingRow.txt}

	\footnotesize \uncover<2->{\hCode{new int[3][]} makes 3 empty places for rows; \hCode{t[0]} is \hCode{null}, ``no row''.
	The message names it: \hCode{"t[0]" is null}. Create every row before using it.}
\end{frame}
% ~~~~~~~~~~~~~~~~~~~~~~~~~~~~~~~~~~~~~~ A frm




% ~~~~~~~~~~~~~~~~~~~~~~~~~~~~~~~~~~~~~~ V frm
\begin{frame}[t,fragile] \frametitle{A ragged table and t[0].length}
	\lstcodedownstep[basicstyle=\ttfamily\scriptsize, firstline=3, lastline=14, firstnumber=3][basicstyle=\ttfamily\fontsize{6.4}{7.6}\selectfont]{codeoutput/RaggedBound.java}{Java}{codeoutput/RaggedBound.txt}

	\footnotesize \uncover<2->{The first loops print too little (row 0 has 1 cell), without an error. The second loops ask
	row 1 for 3 cells and stop. Only \hCode{tri[r].length} is right for every row.}
\end{frame}
% ~~~~~~~~~~~~~~~~~~~~~~~~~~~~~~~~~~~~~~ A frm




% ~~~~~~~~~~~~~~~~~~~~~~~~~~~~~~~~~~~~~~ V
\hSection[Output]{Wrong output, no error}
% ~~~~~~~~~~~~~~~~~~~~~~~~~~~~~~~~~~~~~~ A




% ~~~~~~~~~~~~~~~~~~~~~~~~~~~~~~~~~~~~~~ V frm
\begin{frame}[t,fragile] \frametitle{Printing a table}
	\lstcoderightstep[0.62][basicstyle=\ttfamily\scriptsize]{codeoutput/PrintWrong.java}{Java}{codeoutput/PrintWrong.txt}

	\footnotesize \uncover<2->{Lines 5--6 print only references. A row prints with \hCode{Arrays.toString(row)};
	the whole table with \hCode{Arrays.deepToString(t)}.}
\end{frame}
% ~~~~~~~~~~~~~~~~~~~~~~~~~~~~~~~~~~~~~~ A frm




% ~~~~~~~~~~~~~~~~~~~~~~~~~~~~~~~~~~~~~~ V frm
\begin{frame}[t,fragile] \frametitle{Straight columns with printf}
	\lstcoderightstep[0.62][basicstyle=\ttfamily\scriptsize]{codeoutput/Printf.java}{Java}{codeoutput/Printf.txt}

	\footnotesize \hCode{\%5d}: an \hCode{int} in 5 places, right-aligned; \hCode{\%8.1f}: a \hCode{double} in 8 places
	with one decimal; \hCode{\%n}: a new line. The columns line up because every cell has the same width.
\end{frame}
% ~~~~~~~~~~~~~~~~~~~~~~~~~~~~~~~~~~~~~~ A frm




% ~~~~~~~~~~~~~~~~~~~~~~~~~~~~~~~~~~~~~~ V
\hSection[Debugger]{The debugger}
% ~~~~~~~~~~~~~~~~~~~~~~~~~~~~~~~~~~~~~~ A




% ~~~~~~~~~~~~~~~~~~~~~~~~~~~~~~~~~~~~~~ V frm
\begin{frame}[t,fragile] \frametitle{Watching a table in the debugger}
	\begin{enumerate}\small
		\item Put a breakpoint inside the inner loop of \hCode{Printf} and start \menu{Debug As > Java Application}
		\item In the \menu{Variables} view, \hCode{s} opens to its rows \hCode{[0]}, \hCode{[1]}, \hCode{[2]};
			  each row opens to its cells
		\item Step with \keys{F6}: \hCode{c} changes at every step, \hCode{r} only when a row is done
		\item Add \hCode{s[r][c]} in the \menu{Expressions} view to see the current cell
		\item For \hCode{MissingRow}, the Variables view shows \hCode{t[0]} as \hCode{null} before the crash
	\end{enumerate}

	\vspace{0.5em}
	\small The keys for Windows and macOS are in \hCode{array1D-1-freshman}.
\end{frame}
% ~~~~~~~~~~~~~~~~~~~~~~~~~~~~~~~~~~~~~~ A frm




% ~~~~~~~~~~~~~~~~~~~~~~~~~~~~~~~~~~~~~~ V
\hSection[Summary]{Summary}
% ~~~~~~~~~~~~~~~~~~~~~~~~~~~~~~~~~~~~~~ A




% ~~~~~~~~~~~~~~~~~~~~~~~~~~~~~~~~~~~~~~ V frm
\begin{frame}[t,fragile] \frametitle{Cheat sheet: syntax}


% +++++++++++++++++++++++++++++++++++++++ V lst
%: -Lst. syntax
\begin{lstlisting}[style=codesnippet, language=Java, basicstyle=\ttfamily\scriptsize]
int[][] t = new int[3][4];             // 3 rows of 4, all 0
int[][] m = {{1, 2, 3}, {4, 5, 6}};    // rows in braces
int[][] j = new int[3][];              // rows later; j[0] = new int[5];
t[1][2] = 7;   int x = t[1][2];        // write, read: [row][column]
int[] row = t[1];                      // a whole row
t.length   t[r].length                 // rows; cells of row r
for (int r = 0; r < t.length; r++) {
    for (int c = 0; c < t[r].length; c++) { ... t[r][c] ... }
}
for (int[] r : t) { for (int v : r) { ... } }       // read only
System.out.println(Arrays.deepToString(t));         // print all cells
System.out.printf("%5d", t[r][c]);                  // straight columns
\end{lstlisting}
% +++++++++++++++++++++++++++++++++++++++ A lst


\end{frame}
% ~~~~~~~~~~~~~~~~~~~~~~~~~~~~~~~~~~~~~~ A frm




% ~~~~~~~~~~~~~~~~~~~~~~~~~~~~~~~~~~~~~~ V frm
\begin{frame}[t,fragile] \frametitle{Cheat sheet: errors}


% +++++++++++++++++++++++++++++++++++++++ V tbl
%: -Tbl. errors
\begin{table}
		\small
		\begin{tabular}{@{}L{0.40\textwidth}L{0.52\textwidth}@{}}
			\toprule
			\textbf{Symptom} & \textbf{Cause} \\
			\midrule
			\hCode{']' expected}                 & \hCode{t[1, 2]}, or \hCode{new int[][3]} \\
			\hCode{array required, but int found} & one index too many \\
			\hCode{int[] cannot be converted\dots} & one index too few: \hCode{t[r]} is a row \\
			\hCode{ArrayIndexOutOfBounds}         & rows and columns swapped; \hCode{t[0].length} on a ragged table \\
			\hCode{NullPointerException}          & a row of \hCode{new int[n][]} not created yet \\
			\hCode{[[I@6d06d69c}                  & printed the table itself: use \hCode{Arrays.deepToString} \\
			too few cells printed                & a loop bound from another row \\
			\bottomrule
		\end{tabular}
\end{table}%
% +++++++++++++++++++++++++++++++++++++++ A tbl
\end{frame}
% ~~~~~~~~~~~~~~~~~~~~~~~~~~~~~~~~~~~~~~ A frm




% ~~~~~~~~~~~~~~~~~~~~~~~~~~~~~~~~~~~~~~ V frm
\begin{frame}[t,fragile] \frametitle{Exercises: predict, then run}
	\begin{enumerate}\small
		\item \hCode{ScoresTable}: predict the output, then let the outer loop go to \hCode{QUIZZES}. What happens, and at which line?
		\item \hCode{InitRows}: add \hCode{System.out.println(odd[2][0]);}. Predict the message.
		\item \hCode{LoopKinds}: print the average of every week, using the enhanced for.
		\item \hCode{Printf}: add a header line \hCode{q0 q1 q2 q3 avg} that lines up with the numbers.
		\item \hCode{MissingRow}: create the missing rows with a loop, each of length 4.
		\item Write a program that fills a $5 \times 5$ table with $r \cdot c$ and prints it in straight columns.
	\end{enumerate}
\end{frame}
% ~~~~~~~~~~~~~~~~~~~~~~~~~~~~~~~~~~~~~~ A frm




% ~~~~~~~~~~~~~~~~~~~~~~~~~~~~~~~~~~~~~~ V frm
\begin{frame}[t,fragile] \frametitle{References}
	\small The Java Tutorials, Arrays \cite{TheJavaTutorialsArray}; Big Java \cite{horstmann2016big}.

	\vspace{1em}
	\setbeamertemplate{bibliography item}[text]
	\bibliographystyle{IEEEtran}
	\tiny\bibliography{../../hblib/hblibDataStructures}
\end{frame}
% ~~~~~~~~~~~~~~~~~~~~~~~~~~~~~~~~~~~~~~ A frm




% ~~~~~~~~~~~~~~~~~~~~~~~~~~~~~~~~~~~~~~ V frm
\begin{frame}[t,fragile]
	\begin{center}
		\vspace{4cm}
		\hRemark{\Huge Questions?}
	\end{center}
\end{frame}
% ~~~~~~~~~~~~~~~~~~~~~~~~~~~~~~~~~~~~~~ A frm



\end{document}
